% ===== PRÉAMBULE CANONIQUE — à coller VERBATIM en tête de CHAQUE .tex =====
\documentclass[11pt,a4paper]{article}
\usepackage[utf8]{inputenc}\usepackage[T1]{fontenc}\usepackage{lmodern}
\usepackage[french]{babel}
\usepackage{amsmath,amssymb,mathtools}\usepackage{icomma}
\usepackage{siunitx}\sisetup{locale=FR}  % \qty{}{}, \si{}, \ang{}
\usepackage{xcolor}\usepackage{colortbl}
\usepackage{tikz}\usepackage{pgfplots}\pgfplotsset{compat=1.18}
\usepackage[siunitx,europeanresistors]{circuitikz} % circuits (élec.)
\usepackage[most]{tcolorbox}\usepackage{enumitem}\usepackage{array,booktabs}
\usepackage[a4paper,margin=2.2cm]{geometry}
\usetikzlibrary{arrows.meta,calc,angles,quotes,positioning,patterns,%
  backgrounds,shapes.geometric,decorations.pathmorphing,decorations.markings,intersections}
\definecolor{primary}{RGB}{222,49,99}\definecolor{primarydark}{RGB}{150,22,55}
\definecolor{defcol}{RGB}{0,114,178}\definecolor{methcol}{RGB}{52,92,148}
\definecolor{excol}{RGB}{0,135,90}\definecolor{attncol}{RGB}{200,40,55}
\tcbset{boxbase/.style={enhanced,breakable,boxrule=0.9pt,arc=2.5pt,
  left=3mm,right=3mm,top=2.3mm,bottom=2.3mm,fonttitle=\bfseries,coltitle=white}}
% --- boîtes TUTORIEL (noms EXACTS, ne pas renommer) ---
\newtcolorbox{cle}[1][]{boxbase,colback=defcol!6,colframe=defcol,
  title={\textsc{À retenir}\ifx\relax#1\relax\relax\else\textmd{ : \itshape #1}\fi}}
\newtcolorbox{methode}[1][]{boxbase,colback=methcol!7,colframe=methcol,sharp corners,
  title={\textsc{Méthode}\ifx\relax#1\relax\relax\else\textmd{ --- \itshape #1}\fi}}
\newtcolorbox{exemplebox}[1][]{boxbase,colback=excol!5,colframe=excol,
  title={\textsc{Exemple}\ifx\relax#1\relax\relax\else\textmd{ : \itshape #1}\fi}}
\newtcolorbox{piege}[1][]{boxbase,colback=attncol!6,colframe=attncol,
  title={\textsc{Piège}\ifx\relax#1\relax\relax\else\textmd{ : \itshape #1}\fi}}
\newtcolorbox{remarque}[1][]{boxbase,colback=black!4,colframe=black!45,
  title={\textsc{Remarque}\ifx\relax#1\relax\relax\else\textmd{ : \itshape #1}\fi}}
% --- boîtes EXERCICE / CORRIGÉ (noms EXACTS, auto-comptées) ---
\newtcolorbox[auto counter]{exo}[1][]{boxbase,colback=primary!4,colframe=primary,
  title={\textsc{Exercice}~\thetcbcounter\ifx\relax#1\relax\relax\else\textmd{ --- \itshape #1}\fi}}
\newtcolorbox[auto counter]{corrige}[1][]{boxbase,colback=excol!4,colframe=excol,
  title={\textsc{Corrigé}~\thetcbcounter\ifx\relax#1\relax\relax\else\textmd{ --- \itshape #1}\fi}}
\newcommand{\vv}[1]{\vec{#1}}\newcommand{\ds}{\displaystyle}
\newcommand{\R}{\mathbb{R}}\newcommand{\N}{\mathbb{N}}\newcommand{\Z}{\mathbb{Z}}
% (un \usepackage{fancyhdr}+en-tête est possible ; il est ignoré côté web)
% =========================================================================
\begin{document}
\begin{corrige}[levier équilibré]
\begin{enumerate}[label=\alph*)]
  \item Équilibre de rotation : $G_A\times d_A = G_B\times d_B$.
  \item $G_B = \dfrac{G_A\, d_A}{d_B} = \dfrac{20\times 0{,}3}{0{,}2} = \qty{30}{N}$.
\end{enumerate}
\end{corrige}

\begin{corrige}[à quelle distance ?]
\begin{enumerate}[label=\alph*)]
  \item $d_B = \dfrac{G_A\, d_A}{G_B} = \dfrac{12\times 0{,}5}{30} = \qty{0.2}{m}$.
  \item $d_B=\qty{0.2}{m} < d_A=\qty{0.5}{m}$ : $G_B$ est \emph{plus proche} du pivot. C'est logique :
        le poids le plus grand doit avoir le plus petit bras de levier pour équilibrer les moments.
\end{enumerate}
\end{corrige}

\begin{corrige}[deux moments opposés]
\begin{enumerate}[label=\alph*)]
  \item $\mathcal{M}_1 = F_1\times b_1 = 6\times 0{,}4 = \qty{2.4}{N.m}$ (horaire).
  \item Équilibre : le moment anti-horaire doit égaler le moment horaire :
        $F_2\times b_2 = \mathcal{M}_1$, donc $F_2 = \dfrac{2{,}4}{0{,}3} = \qty{8}{N}$.
\end{enumerate}
\end{corrige}

\begin{corrige}[la balançoire]
\begin{enumerate}[label=\alph*)]
  \item $\mathcal{M} = G\times d = 300\times 1{,}5 = \qty{450}{N.m}$.
  \item Équilibre : $750\times d = 450 \Rightarrow d = \dfrac{450}{750} = \qty{0.6}{m}$. L'adulte, plus
        lourd, s'assoit plus près du pivot.
\end{enumerate}
\end{corrige}

\begin{corrige}[réaction et bras nul]
\begin{enumerate}[label=\alph*)]
  \item La réaction $R_A$ s'applique \emph{en} $A$ : sa ligne d'action passe par l'axe, donc son bras de
        levier est \textbf{nul} ($b=0$).
  \item Son moment est $R_A\times 0 = 0$. Intérêt : en prenant l'axe en $A$, la réaction inconnue $R_A$
        \textbf{disparaît} de l'équation des moments, qui ne contient plus que $R_B$ (et les charges) —
        on l'obtient directement.
\end{enumerate}
\end{corrige}

\end{document}
